\documentclass[12pt,a4paper]{article}
\usepackage[utf8]{inputenc}
\usepackage[T1]{fontenc}
\usepackage[brazil]{babel}
\usepackage{indentfirst}
\usepackage[margin=2.5cm]{geometry}

\title{Resenha Cr\'itica: Artigo 8 ---\\
Managing Technical Debt}
\author{Gustavo Pessoa Firmino Duarte}
\date{19 de Abril de 2026}

\begin{document}
\maketitle

\section{Introdu\c{c}\~ao}

No artigo ``Managing Technical Debt'', Steve McConnell --- fundador e diretor
executivo da Construx Software --- sistematiza um dos conceitos mais importantes
e frequentemente mal compreendidos da engenharia de software moderna: a d\'ivida
t\'ecnica (\textit{technical debt}). O termo foi cunhado por Ward Cunningham em
1992 para descrever a penalidade de desempenho a longo prazo resultante de
decis\~oes de implementa\c{c}\~ao r\'apidas, mas de baixa qualidade. McConnell
amplia e refina essa defini\c{c}\~ao, introduzindo uma taxonomia clara de tipos
de d\'ivida, discutindo o conceito de ``juros'' sobre essa d\'ivida e propondo
um framework pr\'atico para a tomada de decis\~oes sobre quando contrair, manter
ou quitar d\'ivida t\'ecnica.

\section{Tipos de D\'ivida T\'ecnica e o Framework de Decis\~ao}

McConnell classifica a d\'ivida t\'ecnica em dois grandes tipos. O \textbf{Tipo I},
ou d\'ivida n\~ao-intencional, surge de desconhecimento ou descuido: c\'odigo
escrito por um desenvolvedor que simplesmente n\~ao sabia a forma correta de
implement\'a-lo. Esse tipo de d\'ivida \'e involunt\'ario e, em geral, o mais
custoso a longo prazo porque tende a ficar oculto at\'e que seu impacto seja
critic\'o. O \textbf{Tipo II}, ou d\'ivida intencional, \'e uma decis\~ao
deliberada: o time sabe que est\'a tomando um atalho. Ela se subdivide em
\textbf{Tipo II.A} (curto prazo, focalizada), quando h\'a um plano claro de
quitação ap\'os o prazo imediato, e \textbf{Tipo II.B} (longo prazo, n\~ao
focalizada), quando a d\'ivida \'e tomada de forma estrat\'egica com a expectativa
de que ser\'a quitada gradualmente ou nunca.

O autor utiliza a metáfora financeira com rigor: assim como uma d\'ivida monet\'aria
cobra juros, a d\'ivida t\'ecnica cobra ``juros'' na forma de maior esfor\c{c}o
necess\'ario para implementar qualquer funcionalidade futura sobre o c\'odigo
comprometido. Esses juros se acumulam de forma composta: quanto mais c\'odigo \'e
constru\'ido sobre uma base de m\'a qualidade, mais caro fica cada nova altera\c{c}\~ao.

McConnell tamb\'em prop\~oe um framework de tr\^es caminhos para decis\~oes de
desenvolvimento: o \textbf{caminho bom} (Good Path), sem d\'ivida, com tempo
adequado; o \textbf{caminho r\'apido e sujo} (Quick \& Dirty), que contrai d\'ivida
sem plano de quitação; e o \textbf{caminho r\'apido mas n\~ao sujo} (Quick-but-not-dirty),
que encontra solu\c{c}\~oes r\'apidas sem compromenter a qualidade estrutural ---
o caminho ideal quando o prazo \'e imped\'imento real. A visibilidade da d\'ivida
atrav\'es de \textit{debt lists} \'e enfatizada: d\'ivida n\~ao registrada \'e
d\'ivida que nunca ser\'a quitada.

\section{Conex\~ao com a Pr\'atica Profissional}

Este artigo descreve com precis\~ao cir\'urgica uma situa\c{c}\~ao que vivi de
forma intensa: um projeto acad\^emico com prazo de cinco dias usando tecnologias
com as quais eu n\~ao tinha familiaridade pr\'evia. Sob press\~ao, minha estrat\'egia
foi pragm\'atica --- priorizei o ``Make it work'', contando com ferramentas de
IA para acelerar a curva de aprendizado. O resultado foi funcional, mas acumulei
d\'ivida t\'ecnica de Tipo II.A: sabia que os atalhos eram atalhos, mas n\~ao
havia tempo para quet\'a-los no prazo do projeto.

O que o artigo de McConnell me fez perceber \'e que essa d\'ivida n\~ao foi
registrada em lugar algum. N\~ao havia uma \textit{debt list}, nenhum backlog
item marcando os pontos fr\'ageis. Ap\'os a entrega, o projeto foi simplesmente
encerrado --- mas se fosse um produto real, esses ``juros'' se tornariam cada
vez mais pesados nas itera\c{c}\~oes seguintes.

Na ArcelorMittal, observei o peso acumulado de d\'ivida t\'ecnica de Tipo I em
sistemas legados: m\'odulos que ningu\'em na equipe compreendia completamente,
depend\^encias n\~ao documentadas entre sistemas e a resist\^encia institucional
a qualquer mudan\c{c}a por medo de quebrar algo inesperado. Esse cen\'ario \'e
o que McConnell descreve como juros compostos chegando ao limite da insolvência
--- sistemas que consomem praticamente todo o esfor\c{c}o de desenvolvimento
apenas para mant\'e-los funcionando, sem capacidade de evoluir.

\section{Conclus\~ao}

A contribui\c{c}\~ao de McConnell est\'a em transformar d\'ivida t\'ecnica de
uma metáfora vaga em um conceito gerenci\'avel. Ao distinguir tipos de d\'ivida,
ao enfatizar a visibilidade atrav\'es de registros expl\'icitos e ao propor
frameworks de decis\~ao, o autor oferece ferramentas concretas para que equipes
tomem decis\~oes informadas sobre qualidade versus velocidade. A mensagem central
\'e equilibrada: d\'ivida t\'ecnica n\~ao \'e inerentemente m\'a --- \'e uma
ferramenta financeira que, usada conscientemente e com plano de quitação, pode
ser estrat\'egica. Mas d\'ivida acumulada involuntariamente, sem visibilidade e
sem plano, \'e o caminho mais r\'apido para a insolv\^encia t\'ecnica.

\end{document}
